% Esqueleto didático adaptado do pré-projeto fornecido pelo aluno.
\section{Considerações previstas e limitações}

A relevância prática da pesquisa está na possibilidade de apoiar organizações públicas e entidades de interesse público na melhoria do planejamento da compra de passagens aéreas, oferecendo informações mais qualificadas sobre padrões de variação tarifária, faixas de antecedência, riscos de preço e níveis de incerteza por rota. Em vez de substituir integralmente o julgamento administrativo, a abordagem proposta busca fornecer evidências para orientar decisões, qualificar regras internas e apoiar gestores em situações nas quais critérios fixos de antecedência podem ser insuficientes.

No plano científico, a contribuição está em articular Administração Pública, gestão de despesas institucionais e Ciência de Dados aplicada, reconhecendo explicitamente as limitações de inferência causal e contrafactual em bases compostas apenas por compras realizadas. Ao combinar dados internos, dados públicos e coleta prospectiva de cotações, a pesquisa busca desenvolver uma abordagem analítica mais robusta para estudar a relação entre antecedência e variação de preços, sem afirmar indevidamente que o preço observado em uma compra alternativa representaria necessariamente o preço da mesma passagem em outro momento.
